The first operator is positive definite, so the minimum-error strategy always guesses $X=1$; the second has one positive and one negative eigenvalue, so its optimum is a nontrivial $\sigma_z$ measurement.  The observed two-fragment state therefore does not determine the local decoder.\qedhere
\end{proof}

Theorem~\ref{thm:binary-helstrom} is consequently a structured equal-prior result, not a generic binary two-fragment theorem.

\begin{figure*}[t]
\centering
\includegraphics[width=0.98\textwidth]{fig2_binary_illustration.pdf}
\caption{\textbf{Two-fragment binary identifiability and its equal-prior boundary.}
(a) The relevant observable object is built from \emph{same-event} paired fragment data: the two local records must originate from the same physical realization so that the connected correlation $\Omega_{ij}=\rho_{ij}-\rho_i\otimes\rho_j$ is retained.  Pairing records from independent events removes precisely this correlation.
(b) For equal-prior conditionally independent binary records, $\Omega_{ij}$ has operator-Schmidt rank one and its two local factors point along the contrast operators $\Delta_i$ and $\Delta_j$.  Their signs are common outcome relabelings, so the positive/negative eigenspaces determine the local Helstrom readouts without reconstructing $\rho_{i,0}$ and $\rho_{i,1}$ separately.  This is the paper's basic example in which decoder identification is strictly easier than ensemble identification.
(c) The conclusion is not generic once the prior is unequal.  Different latent decompositions can have the same complete two-fragment observation while redistributing the contrast scale between fragments; because the unequal-prior Helstrom operator also depends on that unresolved scale, the optimal local decision boundary can change.  Proposition~\ref{prop:unequal-prior} gives an exact qubit realization of the illustrated ambiguity.}
\label{fig:binary-mechanism}
\end{figure*}

\subsubsection{Residual conditional correlations}

The theorem has a direct robust extension.  Let $\rho_{ij|x}$ be the conditional two-fragment state and define
\begin{equation}
\chi_x=\rho_{ij|x}-\rho_{i|x}\otimes\rho_{j|x}.
\end{equation}
Then
\begin{equation}
\Omega_{ij}=\frac14\Delta_i\otimes\Delta_j+\bar\chi,
\qquad
\bar\chi=\frac{\chi_0+\chi_1}{2}.
\label{eq:robust-connected}
\end{equation}
Let
\begin{equation}
\nu=I(E_i:E_j|X)
\end{equation}
be measured in bits.  Quantum Pinsker and Jensen imply
\begin{equation}
\|\bar\chi\|_2\le\|\bar\chi\|_1
\le\sqrt{2\ln2\,\nu}=: \eta.
\label{eq:cmi-noise}
\end{equation}
The rank-one signal in Eq.~\eqref{eq:robust-connected} has singular value
\begin{equation}
\lambda=\frac14\|\Delta_i\|_2\|\Delta_j\|_2.
\end{equation}

\begin{corollary}[Robust binary decoder direction]
\label{cor:robust-binary}
If $\eta<\lambda$, the leading operator-Schmidt directions $\widehat U_i,\widehat U_j$ of the observed connected operator obey
\begin{equation}
\sin\angle_{\rm HS}(\widehat U_i,\Delta_i)
\le\frac{\eta}{\lambda-\eta},
\label{eq:binary-angle}
\end{equation}
and likewise for $j$, after fixing the free common sign.
\end{corollary}

This isolates the relevant physical ratio: small conditional mutual information is useful only relative to the product of the local record contrasts.  Weak records can be unstable even when the common-cause defect is absolutely small.

\subsubsection{Finite-shot arbitrary-dimensional decoder learning}

A finite-shot theorem can be stated without a free inverse-condition parameter.  Suppose each \emph{same physical event} produces a paired local snapshot $(R_i,R_j)$, with the two local measurement randomizations independent conditional on that event, and
\begin{equation}
\mathbb E[R_i|X=x]=\rho_{i,x},\qquad
\|R_i\|_2\le B_i,
\end{equation}
and similarly for $j$.  The same-event pairing is essential: $\mathbb E[R_i\otimes R_j]$ retains the environmental correlation, whereas snapshots taken from different physical realizations factorize into $\mathbb E R_i\otimes\mathbb E R_j$.  Finite-dimensional classical-shadow constructions are one example of such unbiased snapshot schemes \cite{huang2020}.  The concentration statement below is not dimension-free in an operational sense; the measurement dimension and inverse-channel conditioning enter through $B_i,B_j$.  For example, a product of independent single-qubit Pauli-shadow snapshots on a $q$-qubit fragment has single-shot Hilbert-Schmidt norm $B=5^{q/2}$, so the closed-form bound inherits exponential $5^q$ factors unless additional locality or observable structure is exploited.

From $2N$ same-event realizations, pair consecutive events and define
\begin{equation}
K_s=\frac12(R_i^{2s-1}-R_i^{2s})\otimes(R_j^{2s-1}-R_j^{2s}),
\qquad
\widehat\Omega=\frac1N\sum_{s=1}^{N}K_s.
\end{equation}
Then $\mathbb E\widehat\Omega=\Omega$ and, with probability at least $1-\beta$,
\begin{equation}
\|\widehat\Omega-\Omega\|_2\le
\frac{2B_iB_j}{\sqrt N}\left[1+\sqrt{2\ln(1/\beta)}\right]
=:e_N(\beta).
\label{eq:finite-omega}
\end{equation}

\begin{theorem}[Finite-shot binary readout]
\label{thm:finite-binary}
Under the exact binary product model, if $e_N(\beta)<\lambda$, then with probability at least $1-\beta$,
\begin{equation}
\sin\angle_{\rm HS}(\widehat U_i,\Delta_i)
\le\frac{e_N(\beta)}{\lambda-e_N(\beta)}.
\end{equation}
For target angular error $0<\eps\le1$, it is sufficient to use $n=2N$ physical events with
\begin{equation}
 n\ge128\,
\frac{B_i^2B_j^2}{\eps^2\|\Delta_i\|_2^2\|\Delta_j\|_2^2}
(1+\eps)^2
\left[1+\sqrt{2\ln(1/\beta)}\right]^2.
\label{eq:finite-binary-n}
\end{equation}
With conditional correlations, the same statement holds after replacing $e_N$ by $e_N+\sqrt{2\ln2\,\nu}$.
\end{theorem}

The proof is a Hilbert-space bounded-difference concentration argument followed by Wedin perturbation theory; details are in Appendix~\ref{app:proofs}.  Eq.~\eqref{eq:finite-binary-n} exposes the expected divergence when the local records become indistinguishable.  It should not be read as a universal dimension-independent shadow theorem: the bounded-snapshot constants carry the measurement-interface cost.

\subsection{Multiclass records: subspace, ambiguity, and a third view}
\label{sec:multiclass}

For $r>2$, two fragments retain useful structure but no longer determine the optimal decoder generically.  Let
\begin{equation}
\rho_{ij}=\sum_{x=1}^{r}p_x\rho_{i,x}\otimes\rho_{j,x},
\end{equation}
with
\begin{equation}
\bar\rho_i=\sum_xp_x\rho_{i,x},\qquad
\delta_{i,x}=\rho_{i,x}-\bar\rho_i.
\end{equation}
Then
\begin{equation}
\Omega_{ij}=\rho_{ij}-\bar\rho_i\otimes\bar\rho_j
=\sum_xp_x\delta_{i,x}\otimes\delta_{j,x}.
\label{eq:multi-connected}
\end{equation}
Define the local discriminant operator subspace
\begin{equation}
\cS_i=\Span\{\delta_{i,1},\ldots,\delta_{i,r}\}.
\end{equation}
Since $\sum_xp_x\delta_{i,x}=0$, one has $\dim\cS_i\le r-1$.

\begin{theorem}[Two-view discriminant-subspace identifiability]
\label{thm:subspace}
If $\dim\cS_i=\dim\cS_j=r-1$, then the left and right operator-Schmidt supports of $\Omega_{ij}$ are exactly $\cS_i$ and $\cS_j$, respectively.
\end{theorem}

The subspace is operationally meaningful.  If a Hermitian effect $M=M_{\parallel}+M_{\perp}$ is decomposed with $M_{\perp}\perp\cS_i$, then
\begin{equation}
\tr(M_{\perp}\rho_{i,x})=\tr(M_{\perp}\bar\rho_i)
\end{equation}
for every $x$.  Thus $M_{\perp}$ carries no class-dependent signal.  Pairwise environmental correlations can therefore identify the entire local operator subspace in which a useful multiclass decoder must live, even when they do not identify the individual class states.

\subsubsection{Pairwise data do not determine the multiclass decoder in general}

The remaining coordinates inside $\cS_i$ matter for minimum-error discrimination.  The failure already occurs for commuting states.  Appendix~\ref{app:counterexample} gives an exact three-class construction in which the same two-view diagonal state admits two distinct latent decompositions.  The first has optimal local MAP rule $(0,1,2)$, while the second has $(1,1,2)$.  Because the states commute, these MAP rules are the minimum-error quantum decoders.  Hence
\begin{equation}
\begin{aligned}
&\text{pairwise discriminant-subspace identification}\\[-1mm]
&\hspace{18mm}\not\Rightarrow\ \text{multiclass decoder identification}.
\end{aligned}
\end{equation}

